\section{Intrinsic metrics and varying parameters}\label{sec:intrinsic}

Let $V$ be a finite doubled-quiver representation and suppose
$Z=\mu_{\mathbb C}^{-1}(0)$ contains a module-simple representation
of dimension $m$. Loops and divisible dimension vectors are allowed.
Set $A_\eta=Z\cap\mu^{-1}(\eta)$ and $M_\eta=A_\eta/K$.
An \emph{admissible path} in $M_\eta$ is continuous on a compact
interval with finite subdivision, smooth on each open subinterval
in a representation-type stratum, and of finite one-sided stratum
length at every endpoint. Let $d_\eta$ be the infimum of these
lengths. We include all singular strata.

\begin{theorem}\label{thm:varying-parameters}
Every $(M_\eta,d_\eta)$ is a nonempty proper geodesic space, with
its compact-quotient topology. If $\eta_n\to\eta_*$ and balanced
representatives satisfy $x_n\to x_*$, $y_n\to y_*$, then
\begin{equation}\label{eq:distance-continuity}
 d_{\eta_n}([x_n],[y_n])\longrightarrow
 d_{\eta_*}([x_*],[y_*]).
\end{equation}
Every target point is a limit of balanced representatives at
$\eta_n$, after unitary adjustment. For convergent balanced base
representatives $p_n\to p_*$, the complete pointed spaces
$(M_{\eta_n},d_{\eta_n},[p_n])$ converge to
$(M_{\eta_*},d_{\eta_*},[p_*])$ in pointed Gromov--Hausdorff distance.
\end{theorem}

The upper distance bound needs two separate facts: uniformly short
paths near a singular point and almost isometric paths along strata.

\subsection{Semialgebraic paths and inner distances}
\begin{lemma}\label{lem:uniform-semialgebraic-paths}
For a fixed semialgebraic family
$\mathcal A\subset P\times\mathbb R^N$ of nonempty compact connected
fibers, there is $C$ such that any two points in $\mathcal A_p$
have a continuous semialgebraic joining path of length at most
$C\operatorname{diam}_{\rm Eucl}\mathcal A_p$.
\end{lemma}
\begin{proof}
Semialgebraic triviality partitions the base into finitely many
pieces with fixed compact model fibers. Triangulate each model.
By a finite ordered rule choose simplices containing the endpoints,
join each endpoint to a vertex, and join the vertices along a
spanning tree of the connected one-skeleton. Transport these paths
through the trivializations. The resulting function of parameter,
endpoints and time is semialgebraic and continuous in time; joint
endpoint continuity is unnecessary.
Uniform monotonicity bounds by an integer $J$ the number of
monotonicity intervals of every coordinate of every path. Thus
\[
 \operatorname{Length}(\gamma)
 \le\sum_{\ell=1}^N\operatorname{Var}(\gamma_\ell)
 \le NJ\operatorname{diam}_{\rm Eucl}\mathcal A_p.
\]
The finiteness, triangulation and triviality statements used here
are recalled in \cite[pp.~116, 130, 133, 135]{vdD98}.
\end{proof}

For a closed connected semialgebraic set $A$, its full Euclidean
inner distance $\delta_A$ is finite, since semialgebraic arcs join
its points. It induces the Euclidean topology locally. At $x\in A$,
the component $C$ of $A\cap\overline B_R(x)$ containing $x$ is
compact and connected and contains a relative neighborhood of $x$:
the finitely many other compact components have positive distance
from $x$. Kurdyka--Orro's inner-metric inequality gives
\[
 \delta_C(x,y)\le C_0\|x-y\|^\beta,\qquad 0<\beta\le1
\]
on $C$ \cite[Corollary~2]{KO97}.
Since $\|x-y\|\le\delta_A(x,y)\le\delta_C(x,y)$, the assertion
follows. Closed $\delta_A$-balls are Euclidean bounded and compact:
every Euclidean convergent subsequence converges also in $\delta_A$.

We also require finite-stratum approximations. Apply the near-straight
decomposition of \cite[Proposition~3, pp.~176--178]{KO97} separately
to the finitely many prescribed semialgebraic pieces in a bounded
truncation of $A$. For each $\varepsilon>0$ it yields finitely many
analytic pieces $\Gamma_a$, each in one prescribed piece, whose
closures $F_a$ cover the truncation and admit $C^1$ joining arcs
with interior in $\Gamma_a$ and length at most
$(1+\varepsilon)$ times Euclidean endpoint distance.

For a rectifiable path, start at its first time, choose a member
of the finite closed cover containing that point and visited later,
and jump to its last visit. Repeat until the endpoint. Finiteness
and closedness ensure a choice exists, and no member can recur
after its last visit. The resulting finite chain has total chord
length no greater than the original length. Replace each chord by
its near-straight arc. Smooth each $C^1$ arc inside $\Gamma_a$ on
its open interval, using tolerances tending to zero at its endpoints
and arbitrarily small integrated derivative error.
This preserves endpoints and adds arbitrarily little length.
Finite smooth-stratum paths therefore approximate the full inner
distance, even for an original path with infinitely many crossings.
We do not require geodesics themselves to have finite subdivisions.

\subsection{Fixed singular levels}
The existence of a module-simple representation makes $Z$ reduced, irreducible and
generically module-simple \cite[Theorem~1.2]{CB01}.
A module-simple representation is stable for every parameter,
so Proposition~\ref{prop:real-polystability} gives nonemptiness.
The complex quotient $Q_F$ is irreducible and normal
\cite[Proposition~3.21]{BeSc24}; hence $M_\eta$ is connected by
Corollary~\ref{cor:fixed-face}. The compact quotient
$A_\eta\to M_\eta$ is proper and has connected fibers, since $K$
is connected. Thus $A_\eta$ is connected: a separation of its
domain projects to a separation of the base, because each connected
fiber lies in one part and the quotient is closed.
Properness here also follows by bounding the descended Euclidean
norm over a compact subset of the quotient.

Let $\delta_\eta$ denote the full Euclidean inner metric on $A_\eta$.
The preceding local argument makes it proper and topologically
compatible.
\begin{lemma}\label{lem:quotient-inner-identity}
For $x,y\in A_\eta$,
\begin{equation}\label{eq:quotient-inner-identity}
 d_\eta([x],[y])=\min_{k\in K}\delta_\eta(x,ky).
\end{equation}
The metric $d_\eta$ is proper, complete and geodesic.
\end{lemma}
\begin{proof}
Smooth prequotient representation-type pieces map by Riemannian
submersions to quotient strata
\cite[Theorem~1.1(i), \S2.5]{May22}.
Finite-stratum approximation and projection give the upper bound
for $d_\eta$ in \eqref{eq:quotient-inner-identity}.

Conversely lift each open smooth segment of an admissible quotient
path horizontally, initializing at an interior time.
The lift extends over every compact subinterval: the quotient image
is compact in one smooth stratum and its preimage is compact,
since $K$ is compact and radius is bounded.
The lifted speed equals the quotient speed. Finite endpoint lengths
give Euclidean Cauchy limits in the closed moment level, lying in
the prescribed endpoint orbits. Constant unitary adjustments align
successive segments and the first initial endpoint with $x$.
The concatenation ends at some $ky$ and has exactly the quotient
length, proving the reverse inequality.

The right side is the metric quotient of a proper inner space by
a compact isometric group. It is proper and has the compact-quotient
topology. The subarc inequality and admissible near-minimizers show
it is a length metric. Proper length-space Hopf--Rinow gives
geodesicity and completeness.
\end{proof}
In particular $r([x])=\|x\|$ is $1$-Lipschitz.
At zero level radial paths are horizontal, so
\begin{equation}\label{eq:cone-radius}
 d_0([x],0)=\|x\|.
\end{equation}
This concerns the full quotient, not the completion of its regular
locus alone.

\subsection{Connected neighborhoods in nearby levels}
Pass to a subsequence on which every $\eta_n$ belongs to one face
$F$ of the finite sign arrangement in Section~\ref{sec:slices}.
The face $F_*$ containing $\eta_*$ lies in its closure.
The semistable inclusion gives a projective morphism
\begin{equation}\label{eq:face-map}
 \pi:Q_F\longrightarrow Q_{F_*}
\end{equation}
by \cite[Lemma~2.4]{BeSc24}, with the sign of its character reversed
to match our convention.
The common dense module-simple locus is stable for every character.
Its common geometric quotient is a dense open on which $\pi$ is
an isomorphism. Thus $\pi$ is birational and surjective by properness.
Normality of the integral target gives connected fibers
\cite[Lemma~37.53.6, Tag 0AY8]{Stacks}.
These deductions use the existence of a module-simple point in the complex zero fiber.

Fix $p_*=[x_*]$ and its image $q_*\in Q_{F_*}$.
In a closed algebraic embedding of an affine open chart containing
$q_*$, let $U_s$ be the component of the closed radius-$s$ ball
intersection containing $q_*$. These are compact connected
semialgebraic sets, nested as $s$ decreases, with intersection
$\{q_*\}$. Each contains a relative open neighborhood of $q_*$:
the other finitely many compact components are separated from it.
Their family is semialgebraic by triviality and uniform finiteness
of connected components.

For the fixed algebraic quotient map $q_{F_*}$ on its semistable
open set, put
\begin{equation}\label{eq:balanced-neighborhood}
 B_{\eta,s}=\{x\in A_\eta:q_{F_*}(x)\in U_s\},\qquad\eta\in F.
\end{equation}
The domain contains $A_\eta$, and this is a semialgebraic family
in the original representative space. No inverse balancing map
is assumed semialgebraic or jointly continuous.
Its compact quotient is $\pi^{-1}(U_s)$, which is compact and
connected: a proper surjection with connected fibers has connected
inverse image of a compact connected set, by the closed-map
separation argument. Applying the same argument to connected
$K$-orbits makes $B_{\eta,s}$ compact and connected.

Furthermore
\begin{equation}\label{eq:orbit-collapse}
 \sup_{x\in B_{\eta,s}}\operatorname{dist}_{\rm Eucl}(x,Kx_*)
 \longrightarrow0\quad\text{as }(\eta,s)\to(\eta_*,0),\quad\eta\in F.
\end{equation}
Estimate \eqref{eq:invariant-properness} bounds the representatives,
since their affine invariants range in the compact image of one
$U_{s_0}$. Any contrary sequence has a convergent subsequence with
limit $x\in A_{\eta_*}$. This point is $F_*$-polystable, so the
quotient map is defined near it and gives $q_{F_*}(x)=q_*$.
Balanced-orbit uniqueness forces $x\in Kx_*$, a contradiction.

Taking representatives over any point of $\pi^{-1}(q_*)$ also
proves coverage of the target: they are bounded at $\eta_n$ and
every limiting orbit is $Kx_*$. Unitary adjustment yields
representatives converging to $x_*$.

\begin{lemma}\label{lem:shrinking-diameter}
As $(\eta,s)\to(\eta_*,0)$ with $\eta\in F$,
$\operatorname{diam}_{d_\eta}(B_{\eta,s}/K)\to0$.
\end{lemma}
\begin{proof}
The stabilizer $H=K_{x_*}$ is a connected product of unitary groups.
A small Euclidean normal tube to $Kx_*$ is
$K\times_HN_\varepsilon$, where $N=(T_{x_*}Kx_*)^\perp$,
under $[k,n]\mapsto k(x_*+n)$.
This tube is semialgebraic: its group, orbit, normal space and
action are semialgebraic, and the nearest-point inverse is given
by polynomial quantifiers.
For nearby parameters \eqref{eq:orbit-collapse} puts $B_{\eta,s}$
in the half-sized tube. The slice
$S_{\eta,s}=B_{\eta,s}\cap(x_*+N_\varepsilon)$ is compact and
semialgebraic with $B_{\eta,s}=K\times_HS_{\eta,s}$.
It is connected: connected $H$ preserves each of its finitely
many open-and-closed components, and multiple components would
separate the associated bundle.
Its Euclidean diameter tends to zero by the nearest-point property
and \eqref{eq:orbit-collapse}.
Lemma~\ref{lem:uniform-semialgebraic-paths} gives joining paths
in this fixed family with length tending uniformly to zero.
Projection and \eqref{eq:quotient-inner-identity} give the assertion.
\end{proof}
No connectedness of arbitrary moving moment-level ball intersections
has been assumed.

\subsection{Sharp distance bounds and complete pointed limits}
\begin{proof}[Proof of Theorem~\ref{thm:varying-parameters}]
For the lower bound choose a subsequence with finite liminf of
the distances, and admissible near-minimizers. Lift them from $x_n$
to $k_ny_n$ as in Lemma~\ref{lem:quotient-inner-identity}.
Constant-speed parametrization on $[0,1]$ gives bounded Lipschitz
constants and radius. After subsequences $k_n\to k$ and
Arzel\`a--Ascoli gives a uniform limiting curve in $A_{\eta_*}$
from $x_*$ to $ky_*$. Lower semicontinuity of length gives
\[
 d_{\eta_*}([x_*],[y_*])
 \le\liminf_n d_{\eta_n}([x_n],[y_n]).
\]

For the upper bound take an admissible limit path of length at most
$d_{\eta_*}([x_*],[y_*])+\varepsilon$.
At an interior point of any stratum piece,
Theorem~\ref{thm:metric-slice} identifies the stratum with
$U\times\{0\}$ in \eqref{eq:normal-product}.
Choose the chart radius so that the squared-metric error is less
than $\varepsilon$ at $\eta_*$. Coverage, followed by the common
slice inverse, supplies nearby normal balanced points
$(u_n,z_n)\to0$. Since the moment ignores $U$, $(0,z_n)$ is balanced.

For a compact limit segment $(u(t),0)$ use $(u(t),z_n)$ at
$\eta_n$ and apply the slice map. Its model length is the Euclidean
length of $u$, so comparison on both levels gives, for large $n$,
\[
 \operatorname{Length}(\gamma_n)
 \le\sqrt{\frac{1+2\varepsilon}{1-2\varepsilon}}\,
        \operatorname{Length}(\gamma),\qquad\varepsilon<1/4.
\]
Its stabilizer is independent of $u(t)$, so it lies in one stratum.
Endpoints converge by parameter continuity of the slice and
complementary correction: the bounded retained correction acts
trivially on $U$ as $z_n\to0$.
No weak-eigenvalue inverse or derivative estimate is used.

Around each original junction choose $U_s$ with nearby intrinsic
diameter as small as prescribed by Lemma~\ref{lem:shrinking-diameter}.
Truncate the incident open pieces so their remaining endpoints
are in its relative interior. Continuity and the neighborhood
property permit this; finite endpoint lengths make the removed
total length arbitrarily small. Cover the remaining compact
smooth pieces by finitely many small slice charts and subdivide.
Use the sharp length bound on each segment.

At new subdivision points the two nearby endpoints converge to
the same target point, so their quotient images lie in some
$U_{s_n}$ with $s_n\to0$ by a diagonal choice of nested neighborhoods.
The diameter lemma joins them at cost tending to zero.
At original junctions the truncated endpoints and nearby
approximations lie in the previously chosen $U_s$, giving the
prescribed small cost. The same construction joins to $[x_n]$
and $[y_n]$. Constant or wholly truncated pieces use only these
joins; a constant whole limit path follows from the diameter lemma.
There are finitely many joins, so assign total cost at most
$\varepsilon$. Letting $\varepsilon\to0$ proves the upper bound.
Finite faces and the subsequence principle prove
\eqref{eq:distance-continuity} for the full sequence.

A bounded ball about $[p_n]$ has representatives in a fixed
Euclidean ball because radius is $1$-Lipschitz.
Every sequence in these balls has a subsequence of convergent
representatives, with all limiting distances determined by
\eqref{eq:distance-continuity}. Conversely coverage supplies
nearby representatives for every target point.
Choose a finite $\varepsilon$-net in the compact target ball
of radius $L+1$. Its nearby representatives form an asymptotic
$2\varepsilon$-net of the varying radius-$L$ balls: an uncovered
sequence would have a convergent representative and contradict
distance continuity. Finite distance distortions tend to zero.
Ball inclusions hold with radii $L\pm o(1)$; boundary points can
be approached from inside using geodesics on either side.
This proves pointed convergence on every bounded ball to the
whole proper geodesic target.
\end{proof}

